\section{Experimental Setup}\label{sec:setup}

This section records the hardware and software environment in which every
experiment was executed, the provisions made to keep the study
reproducible, and the ethical considerations that bound its scope.

\subsection{Hardware and Software Stack}\label{ssec:hardware}

All experiments are executed on a single workstation equipped with an
\textbf{NVIDIA GeForce RTX~4060~Ti} graphics processing unit (16~GB GDDR6),
32~GB of system memory, and an AMD Ryzen~7 central processing unit. Using a
single fixed machine for every run keeps the measured latency, GPU peak
memory, training time, and energy figures mutually comparable, since they
are not contaminated by hardware variation across models. The exact Ryzen~7
SKU was not recorded; this does not affect the latency and GPU-memory
figures, which are GPU-bound and CPU-insensitive, but it is relevant to the
wall-clock training time of Table~3, the one CPU-sensitive measurement, and
that table is therefore reported with the explicit caveat given in
Section~\ref{sec:results} that wall-clock training time is not controlled
across runs.

The software stack is held fixed across all runs: Python~3.11 and
PyTorch~2.x with CUDA~12; \texttt{timm}~0.9 for the structural definition of
the donor EfficientNetV2-B0 backbone, of which only the architectural reference
is used (the pre-trained weights are explicitly \emph{not} loaded, in
keeping with the from-scratch training protocol of
Section~\ref{ssec:training}); \texttt{thop} for floating-point operation
counting; \texttt{pytorch-grad-cam}~1.5 for the Grad-CAM++ saliency maps;
\texttt{lime}~0.2 together with \texttt{scikit-image} for the super-pixel
boundary masks; and \texttt{codecarbon} for energy and
CO\textsubscript{2}-equivalent emissions tracking. Deterministic CuDNN is
enabled for every reported run.

\subsection{Reproducibility Provisions}\label{ssec:reproducibility}

The complete experimental package is structured for reproducibility along
five axes.

\begin{enumerate}
  \item \textbf{Code and configuration release.} The full training,
  evaluation, and metric-computation scripts are provided in the
  repository. A single configuration file (\texttt{configs/config.py})
  controls every global hyper-parameter and is enforced across all runs, so
  that the matched recipe of Section~\ref{ssec:training} cannot drift
  between models.
  \item \textbf{Fixed seed.} Every run uses random seed 42 for Python,
  NumPy, and PyTorch, and CuDNN deterministic mode is enabled.
  \item \textbf{Fixed split.} The dataset splits are deterministic functions
  of the seed, which makes train/test contamination across runs impossible
  by construction.
  \item \textbf{Per-run JSON.} Every one of the ten models (the nine
  baseline backbones and the proposed model) emits a
  \texttt{classification\_metrics.json}, a \texttt{performance\_metrics%
  .json}, a \texttt{training\_time.json}, and a per-class textual report
  (\texttt{evaluation\_results.txt}); all of these are committed alongside
  the trained model weights, so that every classification and computational
  number reported in Section~\ref{sec:results} can be independently
  regenerated for both the baseline rows and the proposed model.
  \item \textbf{Explainability artifacts.} The Grad-CAM++ and LIME panels for
  both heads (\texttt{explain\_binary.png} and
  \texttt{explain\_subtype.png}) are produced by a single script
  (\texttt{explain\_model.py}) under identical preprocessing for every
  model.
\end{enumerate}

A small number of runs (the AttentionHub ablation cells) include
\texttt{codecarbon}-measured energy and CO\textsubscript{2} figures, whereas
the older baseline runs report the conservative analytical estimate used in
the absence of \texttt{codecarbon} at the time those runs were executed. We
retain both forms of the measurement, and label them as such, in the
interest of transparency.

\subsection{Ethical Considerations}\label{ssec:ethics}

The datasets used in this study are pre-existing curated collections of oral
lesion images. They are used solely for non-clinical research purposes, and
all results are reported on a held-out split that is disjoint from the
training and validation data. No patient-identifying information is present
in the datasets, and no human subjects are recruited as part of this work.

The proposed model is \emph{not} a clinical decision-support system and
should not be used as one without appropriate clinical validation,
regulatory clearance, and prospective evaluation. We further acknowledge
that automated screening systems can perpetuate disparities present in their
training data: in particular, the demographic composition and
intra-oral-anatomy coverage of the datasets are not characterised in their
published metadata, and external validation across diverse patient
populations is required before any clinical claim can be made. The model is
released under a research-use license, with the explicit intent of enabling
replication and methodological extension rather than direct clinical
deployment.
